\documentclass[../main.tex]{subfiles}

\begin{document}

\section{Trapezoidal Direct Collocation}\label{sec:methodology}

Direct transcription methods discretise the continuous optimal control problem \eqref{eq:OCP} over a finite mesh of nodes and replace the continuous-time dynamics with algebraic relations among the node values~\cite{hargraves1987,betts2010}. The resulting finite-dimensional NLP can be solved by general-purpose nonlinear programming solvers, so I never have to derive the costate equations and transversality conditions of the indirect approach. It is the direct counterpart to the PMP-based formulation used in Homework~1, the same indirect-versus-direct trade-off covered in the course~\cite[Lecture~8]{dclvnotes}.

\subsection{Trapezoidal Quadrature Rule}\label{sec:trap_rule}

The continuous-time dynamics are written compactly as
\begin{equation}
    \dot{\mathbf{x}}(t) = \mathbf{f}\bigl(\mathbf{x}(t),\, \mathbf{u}(t)\bigr),
    \qquad \mathbf{x}\in\mathbb{R}^{n_x},\; \mathbf{u}\in\mathbb{R}^{n_u}.
\end{equation}
Discretise the time horizon $[0, t_f]$ into $N$ uniformly spaced nodes $\{t_k\}_{k=1}^{N}$ with spacing $\Delta t = t_f/(N-1)$, and let
\begin{equation}
    \mathbf{x}_k = \mathbf{x}(t_k),\qquad \mathbf{u}_k = \mathbf{u}(t_k),\qquad
    \mathbf{f}_k = \mathbf{f}(\mathbf{x}_k, \mathbf{u}_k).
\end{equation}
Integrating the state equation over $[t_k, t_{k+1}]$,
\begin{equation}
    \mathbf{x}_{k+1} - \mathbf{x}_k = \int_{t_k}^{t_{k+1}} \mathbf{f}\bigl(\mathbf{x}(\tau),\mathbf{u}(\tau)\bigr)\,\mathrm{d}\tau.
\end{equation}
Approximating the integrand by a linear interpolation between the two endpoints and applying the \textbf{trapezoidal quadrature rule} yields
\begin{equation}
    \mathbf{x}_{k+1} - \mathbf{x}_k \;\approx\; \frac{\Delta t}{2}\bigl(\mathbf{f}_k + \mathbf{f}_{k+1}\bigr).
    \label{eq:trap_rule}
\end{equation}
The order of the rule follows from a Taylor expansion about $t_k$: with $\mathbf{x}_{k+1} = \mathbf{x}_k + \Delta t\,\mathbf{f}_k + \tfrac{1}{2}\Delta t^2\dot{\mathbf{f}}_k + O(\Delta t^3)$ and $\tfrac{\Delta t}{2}(\mathbf{f}_k + \mathbf{f}_{k+1}) = \Delta t\,\mathbf{f}_k + \tfrac{1}{2}\Delta t^2\dot{\mathbf{f}}_k + O(\Delta t^3)$, their difference---the local truncation error---is $\boldsymbol{\tau}_k = -\tfrac{1}{12}\Delta t^3\,\ddot{\mathbf{f}}_k + O(\Delta t^4) = O(\Delta t^3)$; summed over the $O(1/\Delta t)$ intervals this yields a globally second-order, $O(\Delta t^2)$, method~\cite[Eq.~(4.56) and p.~155]{betts2010}. The rule is implicit ($\mathbf{f}_{k+1}$ depends on $\mathbf{x}_{k+1}$) and A-stable: applied to the scalar test equation $\dot{x} = \lambda x$ it gives the amplification factor $x_{k+1}/x_k = (1 + \lambda\Delta t/2)/(1 - \lambda\Delta t/2)$, whose modulus stays below unity for every $\operatorname{Re}(\lambda) < 0$ (the Crank--Nicolson scheme), so stiff modes impose no step-size restriction.

\subsection{Defect Constraints}\label{sec:defects}

Promoting \eqref{eq:trap_rule} from approximation to equality constraint produces the \textbf{defect equations} of trapezoidal direct collocation:
\begin{equation}
    \boldsymbol{\delta}_k \;:=\; \mathbf{x}_{k+1} - \mathbf{x}_k - \frac{\Delta t}{2}\bigl(\mathbf{f}_k + \mathbf{f}_{k+1}\bigr) \;=\; \mathbf{0},
    \qquad k = 1,\, \ldots,\, N-1.
    \label{eq:defects}
\end{equation}
With $n_x$ states this generates $n_x(N-1)$ scalar nonlinear equality constraints. When all defects are satisfied to machine precision, the discrete trajectory matches the continuous-time dynamics to second order in $\Delta t$, and the controls are read as the nodal values of a piecewise-linear input.

For the powered-descent problem ($n_x = 5$, $n_u = 2$) each node contributes $n_x + n_u = 7$ decision variables and the defect block carries $5(N-1)$ rows.

\paragraph{Higher-order option.} The trapezoidal rule is the lowest-order member of the collocation family. Replacing it with the Hermite--Simpson defect
\begin{equation}
    \boldsymbol{\delta}_k = \mathbf{x}_{k+1} - \mathbf{x}_k
    - \frac{\Delta t}{6}\bigl(\mathbf{f}_k + 4\,\bar{\mathbf{f}}_{k} + \mathbf{f}_{k+1}\bigr),
    \qquad
    \bar{\mathbf{x}}_{k} = \tfrac{1}{2}(\mathbf{x}_k+\mathbf{x}_{k+1}) + \tfrac{\Delta t}{8}(\mathbf{f}_k - \mathbf{f}_{k+1}),
    \label{eq:hermite_simpson}
\end{equation}
with $\bar{\mathbf{f}}_{k} = \mathbf{f}(\bar{\mathbf{x}}_{k},\bar{\mathbf{u}}_{k})$ evaluated at the interval midpoint, raises the global order from $O(\Delta t^2)$ to $O(\Delta t^4)$---a three-stage implicit Runge--Kutta (Lobatto~IIIA) scheme~\cite[Eqs.~(4.58)--(4.60), p.~147]{betts2010}. The cost is that the midpoint controls $\bar{\mathbf{u}}_{k}$ become extra NLP unknowns and the path constraints have to be enforced at midpoints as well as nodes. Such a scheme would close the $O(\Delta t^2)$-versus-$O(\Delta t^4)$ accuracy gap observed against the RK4 transcription in Table~\ref{tab:hm2_compare4}.

\subsection{NLP Formulation}\label{sec:nlp}

The decision vector stacks states and controls node by node:
\begin{equation}
    \mathbf{z} = \bigl[\mathbf{x}_1^T,\; \mathbf{u}_1^T,\;
                       \mathbf{x}_2^T,\; \mathbf{u}_2^T,\;
                       \ldots,\;
                       \mathbf{x}_N^T,\; \mathbf{u}_N^T\bigr]^T \in \mathbb{R}^{(n_x+n_u)N}.
    \label{eq:dec_vec}
\end{equation}
The transcribed problem is
\begin{equation}
\begin{aligned}
    \min_{\mathbf{z}} \quad & -m_N \\[2pt]
    \text{s.t.}\quad
    & \mathbf{c}_{\mathrm{eq}}^{\mathrm{lin}}(\mathbf{z}) = \mathbf{0}     && \text{(boundary conditions)}, \\
    & \mathbf{c}_{\mathrm{eq}}^{\mathrm{nl}}(\mathbf{z}) = \mathbf{0}      && \text{(defects)}, \\
    & \mathbf{c}_{\mathrm{ineq}}(\mathbf{z}) \leq \mathbf{0}              && \text{(path constraints)}, \\
    & \mathbf{z}_{\mathrm{lb}} \leq \mathbf{z} \leq \mathbf{z}_{\mathrm{ub}} && \text{(box bounds)}.
\end{aligned}
\label{eq:nlp}
\end{equation}
With $N = 50$, the dimensions are: $7N = 350$ decision variables; $9$ linear equalities (5 initial + 4 final boundary conditions, $m_f$ being free); $5(N-1) = 245$ nonlinear equalities (defects); $4N = 200$ nonlinear inequalities (lower and upper thrust magnitude, two glide-slope half-spaces); and $7N$ box bounds.

\subsection{KKT Optimality Conditions}\label{sec:kkt}

Introducing multipliers $\boldsymbol{\lambda}^{\mathrm{eq}}$ and $\boldsymbol{\mu} \geq \mathbf{0}$ for the equality and inequality constraints, the Lagrangian of \eqref{eq:nlp} is
\begin{equation}
    \mathcal{L}(\mathbf{z},\boldsymbol{\lambda}^{\mathrm{eq}},\boldsymbol{\mu})
    = J(\mathbf{z}) + (\boldsymbol{\lambda}^{\mathrm{eq}})^T \mathbf{c}_{\mathrm{eq}}(\mathbf{z})
                    + \boldsymbol{\mu}^T \mathbf{c}_{\mathrm{ineq}}(\mathbf{z}).
\end{equation}
Under a constraint qualification such as the linear independence of the active constraint gradients (LICQ), a necessary condition for $\mathbf{z}^*$ to be a local minimiser is the existence of multipliers $(\boldsymbol{\lambda}^{\mathrm{eq},*},\boldsymbol{\mu}^*)$ satisfying the \textbf{Karush--Kuhn--Tucker} (KKT) conditions~\cite[Ch.~12]{nocedal2006}:
\begin{align}
    \nabla_{\mathbf{z}} \mathcal{L}(\mathbf{z}^*,\boldsymbol{\lambda}^{\mathrm{eq},*},\boldsymbol{\mu}^*) &= \mathbf{0}
        && \text{(stationarity),}\\
    \mathbf{c}_{\mathrm{eq}}(\mathbf{z}^*) = \mathbf{0},\quad \mathbf{c}_{\mathrm{ineq}}(\mathbf{z}^*) &\leq \mathbf{0}
        && \text{(primal feasibility),}\\
    \boldsymbol{\mu}^* &\geq \mathbf{0}
        && \text{(dual feasibility),}\\
    \mu_j^*\, c_{\mathrm{ineq},j}(\mathbf{z}^*) &= 0,\quad \forall j
        && \text{(complementary slackness).}
\end{align}
Complementarity says that an inequality is either active ($c_j = 0$, with $\mu_j > 0$) or strictly inactive ($c_j < 0$, with $\mu_j = 0$). At the optimum of the powered-descent problem the upper thrust bound $\|\mathbf{T}\| = T_{\max}$ is active on the two burn arcs (about half of the nodes), which matches the max--coast--max structure of the fuel-optimal control law for linear-in-control dynamics. The corresponding KKT multipliers are positive, so they flag the binding constraint that shapes the solution.

\subsection{Fuel-Optimal Thrust Structure}\label{sec:bangbang}

Even though the problem is solved directly, its optimal thrust profile carries the structure that the maximum principle predicts for dynamics linear in the control. Forming the Hamiltonian of the continuous OCP~\eqref{eq:OCP},
\begin{equation}
    \mathcal{H} = \lambda_x v_x + \lambda_y v_y
    + \frac{\lambda_{v_x}T_x + \lambda_{v_y}T_y}{m}
    - \lambda_{v_y}\, g - \lambda_m\frac{\|\mathbf{T}\|}{c},
    \label{eq:hm2_hamiltonian}
\end{equation}
and aligning the thrust direction with the primer vector $\mathbf{p} = (\lambda_{v_x}, \lambda_{v_y})/m$, the Hamiltonian becomes linear in the thrust magnitude $\|\mathbf{T}\| \in [T_{\min}, T_{\max}]$, with coefficient the switching function
\begin{equation}
    S = \|\mathbf{p}\| - \frac{\lambda_m}{c}.
    \label{eq:hm2_switching}
\end{equation}
The optimal magnitude therefore lies at a bound---$T_{\max}$ or $T_{\min}$ according to the sign of $S$, with a singular arc only if $S \equiv 0$ on an interval---producing the \textbf{max--coast--max} (bang--off--bang) profile~\cite[Sec.~3.9]{bryson1975}. This is the same bang--bang mechanism behind the bilinear-tangent law of Homework~1; here it shows up numerically as the active set of the upper thrust bound (Section~\ref{sec:kkt}) and is visible in the thrust histories of Sections~\ref{sec:results} and~\ref{sec:task2}.

\subsection{Discrete Multipliers as Costate Estimates}\label{sec:covector}

The NLP multipliers are not just algebraic by-products: as the mesh is refined they converge to scaled values of the continuous costates, the covector mapping property. For the trapezoidal transcription, the discrete multipliers $\boldsymbol{\alpha}_k$ of the path constraints map to the continuous path-constraint adjoint $\hat{\boldsymbol{\mu}}_k$ (distinct from the inequality multipliers $\boldsymbol{\mu}$ of Section~\ref{sec:kkt}) through~\cite[Sec.~4.11.2, Eqs.~(4.235)--(4.237)]{betts2010}
\begin{equation}
    \hat{\boldsymbol{\mu}}_1 = -\frac{2}{h_1}\,\boldsymbol{\alpha}_1,\qquad
    \hat{\boldsymbol{\mu}}_k = -\frac{2}{h_{k-1}+h_k}\,\boldsymbol{\alpha}_k,\qquad
    \hat{\boldsymbol{\mu}}_M = -\frac{2}{h_{M-1}}\,\boldsymbol{\alpha}_M,
    \label{eq:covector}
\end{equation}
with $h_k$ the local step. This is why the converged active set is physically meaningful: the nodes where the thrust-bound multiplier is non-zero are the powered (constrained) arcs of the continuous solution, so the NLP recovers the max--coast--max structure of Section~\ref{sec:bangbang} without ever forming the costates explicitly.

\subsection{Sequential Quadratic Programming}\label{sec:sqp}

The NLP \eqref{eq:nlp} is solved with MATLAB's \texttt{fmincon} configured with the SQP algorithm~\cite[Ch.~18]{nocedal2006}. SQP iteratively constructs a quadratic-programming approximation of the KKT system: at each iterate $\mathbf{z}^{(k)}$ the search direction $\mathbf{d}^{(k)}$ is found by solving
\begin{equation}
\begin{aligned}
    \min_{\mathbf{d}} \quad & \tfrac{1}{2}\mathbf{d}^T \mathbf{H}^{(k)} \mathbf{d} + \nabla J(\mathbf{z}^{(k)})^T \mathbf{d} \\[2pt]
    \text{s.t.}\quad
    & \nabla \mathbf{c}_{\mathrm{eq}}(\mathbf{z}^{(k)})^T \mathbf{d} + \mathbf{c}_{\mathrm{eq}}(\mathbf{z}^{(k)}) = \mathbf{0}, \\
    & \nabla \mathbf{c}_{\mathrm{ineq}}(\mathbf{z}^{(k)})^T \mathbf{d} + \mathbf{c}_{\mathrm{ineq}}(\mathbf{z}^{(k)}) \leq \mathbf{0},
\end{aligned}
\end{equation}
where $\mathbf{H}^{(k)}$ is a quasi-Newton (BFGS) approximation of the Hessian of the Lagrangian. The next iterate is $\mathbf{z}^{(k+1)} = \mathbf{z}^{(k)} + \alpha^{(k)} \mathbf{d}^{(k)}$, with the step length $\alpha^{(k)}$ obtained from a line search on a merit function that balances objective decrease against constraint violation.

The Jacobian of the defect block~\eqref{eq:defects} is highly sparse: each defect $\boldsymbol{\delta}_k$ couples only the two adjacent nodes $k$ and $k+1$, through constant $\pm\mathbf{I}$ blocks plus the nodal values of $\mathbf{f}$, giving a block-bidiagonal structure. Betts uses this sparsity, grouping the finite-difference perturbations so that the cost of building the Jacobian scales with the bandwidth rather than the full variable count~\cite[Sec.~4.6]{betts2010}. Here \texttt{fmincon} estimates a dense finite-difference Jacobian, which is fine at $N = 50$ (350 variables) but is the first thing I would replace---by supplying the sparsity pattern---to scale the mesh up.

\end{document}
